\scanpage{97}
% archive sheet number: 96
\uline{Corollaire} 3.4. Soient $S$ un préschéma, $u:G\quad H$ un homomorphisme
de \deleted{$S$-groupes} $S$-préschémas en groupes localement diagonalisables\inserted{ avec $H$ de type fini sur $S$}, Alors:

a) $G'=\operatorname{Ker}u$ est localement diagonalisable, il est de type fini sur $S$ si $G$ l'est.

b) Le quotient $G/G'$ existe, de façon plus précise la relation
d'équivalence définie par $G'$ dans $G$ est effective pour la descente
fidèlement plate quasi-compacte (cf exposé IV). De plus $G/G\deleted{\unsure{x|*}}'$ est
localement diagonalisable.

c) \deleted{Par suite, (loc.cit)} L'homomorphisme $u:G\quad H$ se factorise
de façon unique en $G\overset v{\quad}G/G'\overset w{\quad}H$, où $v$ est l'homomorphisme
canonique (donc $v$ est fidèlement plat et quasi-compact)\deleted{\unsure{, et|,et}} De plus
$w$ est une immersion fermée, et à fortiori un monomorphisme.

La première assertion de c) est conséquence de b), par définition
du quotient $G/G'$ . Comme toutes les assertions faites sont de nature
locale sur $S$, on peut alors supposer $G$ et $H$ \deleted{définis par} diagonalisables , \deleted{définis} de la forme $D_S(M)$ et $D_{\unsure{3|S}}(N)$, et (compte tenu de 1.5.
et 2.1. b) \inserted{qui implique $N$ de type fini}) $u$ défini par un homomorphisme $u':N\quad M$. Alors en
vertu de 3.1. $G'$ est isomorphe à $D_S(\operatorname{Coker}u')$, ce qui prouve a). Pour
prouver b) et c) , il suffit de réaliser $G/G'$ comme $D_S(\operatorname{Im}u')$, en
appliquant 3.3.

\uline{Remarques} Le résultat d'existence de quotients 3.4. sera substantiellement généralisé au N° suivant, \deleted{\unsure{ainsi que}} On notera que dans le présent
N° et le précédent, l'hypothèse $S\ne\varnothing$ n'est intervenue que pour assurer la validité de certaines réciproques, \deleted{affirmant que si $G=D_S$} permettant de certaines hypothèses sur des schémas en groupes sur $S$, de
conclure des propriétés correspondantes pour certains groupes ordinaires
Les résultats en sens inverse sont valables sans restriction sur $S$, et
les démonstrations données ici s'appliquent dans le cas général.
